\subsection{The case of groups without arbitrarily small finite subgroups}

THEOREM 2. --- Let G be a locally compact group satisfying the following condition:

\begin{enumerate}
\def\labelenumi{(\Alph{enumi})}
\setcounter{enumi}{11}
\tightlist
\item
  There exists a neighborhood of \(e\) in \(G\) that contains no finite subgroup of \(G\) not reduced to \(e\) .
\end{enumerate}

The following properties then hold:

\begin{enumerate}
\def\labelenumi{(\roman{enumi})}
\item
  The set D of discrete subgroups of G is locally closed in \(\Sigma\) (which is equivalent to saying that it is locally compact).
\item
  For a closed subset \(\mathbf{A}\) of \(\mathbf{D}\) to be compact, it is necessary and sufficient that there exist a neighborhood \(\mathbf{U}\) of \(e\) in \(\mathbf{G}\) such that \(\mathbf{H} \cap \mathbf{U} = \{e\}\) for every subgroup \(\mathbf{H} \in \mathbf{A}\) .
\item
  If in addition G is generated by a compact neighborhood of e, then the set \(D_{c}\) of discrete subgroups H of G such that G/H is compact is locally closed in \(\Sigma\) (hence is locally compact).
\end{enumerate}

We have \(\mathbf{D}_c = \mathbf{D} \cap \Sigma_c^0\) , therefore (iii) is a consequence of (i), and Th. 1 (ii) of No.~3.

With the notations of No.~3, Prop. 7, it suffices, for proving (i) and (ii), to prove that:

PROPOSITION 8. --- The canonical bijection of N onto D is a homeomorphism.

Now, if \(\Gamma_d\) is the set of Haar measures on the discrete subgroups of G, then D is canonically homeomorphic to the space of orbits of the group of homotheties in \(\Gamma_d\) with ratio \(>0\) (GT, I, §5, No.~2, Prop. 4). It therefore suffices to prove that the canonical mapping \(\alpha \mapsto (\alpha(\{e\}), \alpha/\alpha(\{e\}))\) of \(\Gamma_d\) onto \(\mathbf{R}_+^* \times \mathbf{N}\) is a homeomorphism, which will result from the following lemma:

Lemma 4. --- If the group G satisfies the condition (L), the mapping \(\alpha \mapsto \alpha (\{e\})\) of \(\Gamma_d\) into \(\mathbf{R}_+^*\) is vaguely continuous.

Let us consider a measure \(\alpha \in \Gamma_d\) ; let \(V_0\) be a relatively compact open neighborhood of \(e\) in \(G\) such that \(H_\alpha \cap V_0 = \{e\}\) and such that there exists no finite subgroup of \(G\) contained in \(V_0\) and not reduced to \(e\) . Let \(V\) be a symmetric compact neighborhood of \(e\) such that \(V^3 \subset V_0\) , and let \(U\) be a symmetric neighborhood of \(e\) such that \(U^2 \subset V\) . Let \(\varphi\) (resp. \(\psi\) ) be a function in \(\mathcal{K}_+(G)\) , with values in \([0,1]\) , equal to 1 on \(V^3\) (resp. at the point \(e\) ) and with support contained in \(V_0\) (resp. in \(U\) ). The set of measures \(\beta \in \Gamma_d\) such that \(|\beta(\varphi) - \alpha(\varphi)| \leqslant \varepsilon\) and \(|\beta(\psi) - \alpha(\psi)| \leqslant \varepsilon\) is a neighborhood \(W\) of \(\alpha\) . We propose to show that, provided \(\varepsilon\) is taken to be sufficiently small, \(H_\beta \cap V = \{e\}\) for every \(\beta \in W\) ; it will then follow that \(\beta(\psi) = \beta(\{e\})\) , hence that \(|\beta(\{e\}) - \alpha(\{e\})| \leqslant \varepsilon\) , which will prove the lemma.

It will suffice to show that, for \(\beta \in \mathbf{W}\) ,

\[
(\mathrm{V} ^ {2} - \mathrm{V}) \cap \mathrm{H} _ {\beta} = \varnothing .\tag{8}
\]

For, suppose that this point is established: then, for \(x\) and \(y\) in \(\mathbf{V} \cap \mathbf{H}_{\beta}\) , one has \(xy^{-1} \in \mathbf{V}^2 \cap \mathbf{H}_{\beta}\) ; but, by virtue of (8), this implies \(xy^{-1} \in \mathbf{V} \cap \mathbf{H}_{\beta}\) ; in other words, \(\mathbf{V} \cap \mathbf{H}_{\beta}\) is a subgroup of \(\mathbf{G}\) , which is obviously discrete and compact, hence finite; but then, by the choice of \(\mathbf{V}_0\) , this implies that indeed \(\mathbf{V} \cap \mathbf{H}_{\beta} = \{e\}\) .

Let us argue by contradiction and so assume that there exists a point \(z\) of \(\mathbf{V}^2 -\mathbf{V}\) that belongs to \(\mathbf{H}_{\beta}\) ; by the choice of U and V, we have \(\psi (sz^{-1}) + \psi (s)\leqslant \varphi (s)\) in G, the relation \(z\notin \mathrm{U}^2\) implying \(\mathrm{U}z\cap \mathrm{U} = \emptyset\) . Since

\[
\int \psi (s z ^ {- 1}) d \beta (s) = \int \psi (s) d \beta (s),
\]

it follows that \(2\beta (\psi)\leqslant \beta (\varphi)\leqslant \alpha (\varphi) + \varepsilon\) ; but we also have

\[
\beta (\psi) \geqslant \alpha (\psi) - \varepsilon ,
\]

and by construction \(\alpha(\varphi) = \alpha(\psi) = \alpha(\{e\})\) . We thus arrive at a contradiction by taking \(\varepsilon < \alpha(\{e\})/3\) . Q.E.D.

In graphic terms, a group G satisfying the condition (L) is said to not have arbitrarily small finite subgroups. *It can be shown that every Lie group satisfies the condition (L); but this condition is not characteristic of Lie groups; for example, the multiplicative group of \(p\) -adic integers congruent to 1 mod \(p\) satisfies (L).*

